\subsection{Krull domains}

DEFINITION 3. An integral domain A is called a Krull domain if there exists a family

\((v_{\iota})_{\iota\in I}\) of valuations on the field of fractions \(\mathbf{K}\) of A with the following properties:

(AK \(_{I}\) ) the valuations \(v_{\iota}\) are discrete;

(AK \(_{II}\) ) the intersection of the rings of the v \(_{\iota}\) is A;

\((\mathbf{A}\mathbf{K}_{\mathrm{III}})\) for all \(x \in \mathbf{K}^*\) , the set of indices \(\iota \in \mathbf{I}\) such that \(v_{\iota}(x) \neq 0\) is finite.

It obviously suffices to verify condition \((\mathrm{AK}_{\mathrm{III}})\) for the elements x of A - (0).

\textbf{Examples}\par

\begin{enumerate}
\def\labelenumi{(\arabic{enumi})}
\item
  Every discrete valuation ring is a Krull domain.
\item
  More generally, every principal ideal domain A is a Krull domain. For let \((p_{\iota})_{\iota\in I}\) be a representative system of extremal elements of A and let \(v_{\iota}\) be the valuation on the field of fractions of A defined by \(p_{\iota}\) (Chapter VI, § 3, no. 3, Example 4). It is immediately seen that the family \((v_{\iota})_{\iota\in I}\) satisfies properties \((\mathbf{AK}_{\mathbf{I}}), (\mathbf{AK}_{\mathbf{II}})\) and \((\mathbf{AK}_{\mathbf{III}})\) .
\item
  Let F be a field and \((\mathbf{R}_{i})_{1\leqslant i\leqslant n}\) a finite family of subrings of F which are Krull domains. Then their intersection \(S=\bigcap_{j=1}^{n}R_{j}\) is a Krull domain. For \(1\leqslant j\leqslant n\) let \((v_{j\iota})_{\iota\in I_j}\) be a family of valuations on the field of fractions of \(R_{j}\) satisfying \((\mathrm{AK}_{\mathrm{I}}),(\mathrm{AK}_{\mathrm{II}}),(\mathrm{AK}_{\mathrm{III}})\) (where A is replaced by \(R_{j}\) ). Let \(w_{j\iota}\) denote the restriction of \(v_{j\iota}\) to the field of fractions of S. Then the family \((v_{j\iota})_{1\leqslant j\leqslant n,\iota\in I_j}\) obviously satisfies \((\mathrm{AK}_{\mathrm{II}})\) (where A is replaced by S) and also \((\mathrm{AK}_{\mathrm{III}})\) since the set of indices j is finite. The valuations \(w_{j\iota}\) are either discrete or improper. By retaining only those which are discrete, a family is obtained which obviously satisfies \((\mathrm{AK}_{\mathrm{I}}),(\mathrm{AK}_{\mathrm{II}})\) and \((\mathrm{AK}_{\mathrm{III}})\) (where A is replaced by S). Hence S is certainly a Krull domain.
\item
  In particular, if A is a Krull domain and \(K'\) a subfield of the field of fractions K of A, \(K' \cap A\) is a Krull domain.
\end{enumerate}

THEOREM 2. Let A be an integral domain. For A to be a Krull domain, it is necessary and sufficient that the two following conditions be satisfied:

\begin{enumerate}
\def\labelenumi{(\alph{enumi})}
\item
  A is completely integrally closed;
\item
  every non-empty family of divisorial integral ideals of \(\mathbf{A}\) admits a maximal element (with respect to the relation \(\subset\) ).
\end{enumerate}

Moreover, if \(\mathrm{P}(\mathrm{A})\) is the set of extremal elements of \(\mathrm{D}(\mathrm{A})\) , \(\mathrm{P}(\mathrm{A})\) is then a basis of the Z-module \(\mathrm{D}(\mathrm{A})\) and the positive elements of \(\mathrm{D}(\mathrm{A})\) are the linear combinations of the elements of \(\mathrm{P}(\mathrm{A})\) with coefficients \(\geqslant0\) .

Let A be a Krull domain. It is completely integrally closed (Chapter VI, § 4, no. 5, Corollary to Proposition 9). Let \((v_{\iota})_{\iota\in I}\) be a family of valuations on the field of fractions K of A satisfying \((\mathrm{AK}_{\mathrm{I}})\) , \((\mathrm{AK}_{\mathrm{II}})\) and \((\mathrm{AK}_{\mathrm{III}})\) . The \(v_{\iota}\) may be assumed to be normed (Chapter VI, § 3, no. 6, Definition 3). For all \(\mathfrak{a} \in \mathrm{I}(\mathbf{A})\) , we shall write:

\[
v _ {\iota} (\mathfrak{a}) = \sup _ {\mathfrak{a} \subset Ax} (v _ {\iota} (x));\tag{3}
\]

then \(v_{\iota}(\mathfrak{a}) \in \mathbf{Z}\) , for, if \(a\) is a non-zero element of \(\mathfrak{a}\) , the relation \(Ax \supset Aa\) implies that \(v_{\iota}(x) \leqslant v_{\iota}(a)\) (by \((\mathrm{AK}_{\mathrm{II}})\) ), which shows that the family of \(v_{\iota}(x)\) ( \(\mathfrak{a} \subset Ax\) ) is bounded above. We establish the following properties:

\begin{enumerate}
\def\labelenumi{(\arabic{enumi})}
\tightlist
\item
  Let \(\mathfrak{a}\) be a divisorial fractional ideal; in order that \(y \in \mathfrak{a}\) , it is necessary and sufficient that \(v_{\iota}(y) \geqslant v_{\iota}(\mathfrak{a})\) for all \(\iota \in \mathbf{I}\) .
\end{enumerate}

As \(\mathfrak{a}\) is divisorial, the relation \(y \in \mathfrak{a}\) is equivalent to the relation `` \(\mathfrak{a} \subset \mathbf{A}x\) implies \(y \in \mathbf{A}x\) ''. Now, by \((\mathrm{AK}_{\mathrm{II}})\) , the relation \(y \in \mathbf{A}x\) is equivalent to `` \(v_{\iota}(y) \geqslant v_{\iota}(x)\) for all \(\iota \in \mathbf{I}\) ''. Whence (1).

\begin{enumerate}
\def\labelenumi{(\arabic{enumi})}
\setcounter{enumi}{1}
\tightlist
\item
  Let \(\mathfrak{a}\) and \(\mathfrak{b}\) be two divisorial fractional ideals of \(\mathbf{A}\) ; in order that \(\mathfrak{a} \subset \mathfrak{b}\) , it is necessary and sufficient that \(v_{\iota}(\mathfrak{a}) \geqslant v_{\iota}(\mathfrak{b})\) for all \(\iota \in \mathbf{I}\) .
\end{enumerate}

This follows immediately from property (1).

\begin{enumerate}
\def\labelenumi{(\arabic{enumi})}
\setcounter{enumi}{2}
\tightlist
\item
  If \(x \in \mathbf{K}^*\) , then \(v_{\iota}(\mathbf{A}x) = v_{\iota}(x)\) .
\end{enumerate}

If \(A y \supset A x\) , then \(v_{\iota}(y) \leqslant v_{\iota}(x)\) by \((\mathbf{A K}_{\mathrm{II}})\) and the minimum value of \(v_{\iota}(y)\) is taken at \(y = x\) .

\begin{enumerate}
\def\labelenumi{(\arabic{enumi})}
\setcounter{enumi}{3}
\tightlist
\item
  For all \(\mathfrak{a} \in \mathbf{I}(\mathbf{A})\) , the indices \(\iota \in \mathbf{I}\) such that \(v_{\iota}(\mathfrak{a}) \neq 0\) are finite in number.
\end{enumerate}

There exist \(x, y\) in \(\mathbf{K}^*\) such that \(Ax \subset \mathfrak{a} \subset Ay\) . By properties (2) and (3), \(v_{\iota}(x) \geqslant v_{\iota}(\mathfrak{a}) \geqslant v_{\iota}(y)\) for all \(\iota \in I\) . It then suffices to apply \((\mathrm{AK}_{\mathrm{III}})\) .

We have therefore shown the following lemma:

LEMMA 1. If A is a Krull domain and \((v_{\iota})_{\iota \in I}\) is a family of normed valuations on K satisfying \((\mathrm{AK}_{\mathrm{I}})\) , \((\mathrm{AK}_{\mathrm{II}})\) and \((\mathrm{AK}_{\mathrm{III}})\) , the mapping \(\mathfrak{a} \mapsto (v_{\iota}(\mathfrak{a}))_{\iota \in I}\) is a decreasing injective mapping of the set of divisorial integer ideals of A (ordered by \(\subset\) ) to the set of positive elements of the ordered group the direct sum \(\mathbf{Z}^{(I)}\) .

This being so, every non-empty set of positive elements of \(\mathbf{Z}^{(I)}\) has a minimal element (Algebra, Chapter VI, § 1, no. 13, Theorem 2). Hence A certainly satisfies property (b) of the statement.

Conversely, let A be an integral domain satisfying properties (a) and (b) of the statement. Since A is completely integrally closed, D(A) is an ordered group (no. 2, Theorem 1). This group is a lattice (no. 1, Proposition 2). By condition (b) of the statement, every non-empty family of positive elements of D(A) has a minimal element. Let P(A) be the set of extremal elements of D(A). Then (Algebra, Chapter VI, § 1, no. 13, Theorem 2) P(A) is a basis of the Z-module D(A) and the positive elements of D(A) are the linear combinations with positive integer coefficients of the elements of P(A).
 
Thus, for \(x \in \mathbf{K}^*\) , rational integers \(v_{\mathbb{P}}(x)\) are defined (for \(\mathbf{P} \in \mathbb{P}(\mathbf{A})\) ) by writing:

\[
\operatorname{div} (x) = \sum_ {\mathbf {P} \in \mathbf {P} (\mathbf {A})} v _ {\mathbf {P}} (x) \cdot \mathbf {P}.\tag{4}
\]

We also write \(v_{\mathbf{P}}(0) = +\infty\) .

From the relations

\[
\operatorname{div} (x y) = \operatorname{div} (x) + \operatorname{div} (y)
\]

and

\[
\operatorname{div} (x + y) \geqslant \inf (\operatorname{div} (x), \operatorname{div} (y)),
\]

for \(x, y\) and \(x + y\) in \(\mathbf{K}^*\) , we deduce that the \(v_{\mathbb{P}}\) are discrete valuations on \(\mathbf{K}\) . In order that \(x \in \mathbf{A}\) , it is necessary and sufficient that \(\mathrm{div}(x) \geqslant 0\) , that is that \(v_{\mathbb{P}}(x) \geqslant 0\) for all \(\mathbb{P} \in \mathbb{P}(\mathbf{A})\) . Thus the \(v_{\mathbb{P}}\) satisfy conditions \((\mathbf{AK}_{\mathbf{I}})\) and \((\mathbf{AK}_{\mathbf{II}})\) and obviously also \((\mathbf{AK}_{\mathbf{III}})\) .

COROLLARY. For a Noetherian ring to be a Krull domain, it is necessary and sufficient that it be an integrally closed domain.

An integrally closed Noetherian domain is completely integrally closed (Chapter V, § 1, no. 4).

There are non-Noetherian Krull domains, for example the polynomial ring \(K[X_{n}]_{n\in\mathbf{N}}\) over a field K in an infinity of indeterminates (cf.~Exercise 8).

